\chapter*{}

\section{Ruben}

\subsection {Cantrips}
\DndSpellHeader%
  {\textbf{Toll the Dead}}
  {Necromancy Cantrip}
  {1 action}
  {60 feet}
  {V, S}
  {Instantaneous}
  %{\textbf{Attack/Save}: WIS 13}\\\
  
You point at one creature you can see within range, and the sound of a dolorous bell fills the air around it for a moment. The target must succeed on a Wisdom saving throw or take 1d8 necrotic damage. If the target is missing any of its hit points, it instead takes (1d12) necrotic damage.

{\textbf{Cantrip Upgrade.}} The spell’s damage increases by one die when you reach 5th level (2d8 or 2d12), 11th level (3d8 or 3d12), and 17th level (4d8 or 4d12).


\DndSpellHeader%
{\textbf{Vengeful Blade}}
{Evocation Cantrip}
{1 action}
{Self}
{S, M (a melee weapon worth at least 1 gp)}
{Instantaneous}
%{\textbf{Attack/Save}: WIS 13}\\\

You brandish the weapon used in the spell’s casting and make a melee attack with it against one creature within 5 feet of you. On a hit, the target suffers the weapon attack’s normal effects and then radiates a dark aura of energy until the start of your next turn. If the target makes an attack or casts a spell before then, the target takes (1d8) necrotic damage and the spell ends.

\textbf{Cantrip Upgrade.} This spell’s damage increases when you reach certain levels. At 5th level, the melee attack deals an extra 1d8 necrotic damage to the target on a hit, and the damage the target takes for making an attack or casting a spell increases to 2d8. Both damage rolls increase by 1d8 at 11th level (2d8 and 3d8) and again at 17th level (3d8 and 4d8).


\subsection {Class Features}
\begingroup
\DndSetThemeColor[PhbMauve]
\begin{DndComment}{Second Wind}
	You gain a Fighting Style feat of your choice, and whenever you gain a Fighter level, you can replace the feat you chose with a different Fighting Style feat.
\end{DndComment}

\begin{DndComment}{Spellcasting}
	You have learned to cast spells through prayer and meditation.
	
	Spell Slots. The Paladin Features table shows how many spell slots you have to cast your level 1+ spells. You regain all expended slots when you finish a Long Rest.
	
	Prepared Spells of Level 1+. You prepare the list of level 1+ spells that are available for you to cast with this feature. To start, choose two level 1 Paladin spells.
	
	The number of spells on your list increases as you gain Paladin levels, as shown in the Prepared Spells column of the Paladin Features table. Whenever that number increases, choose additional Paladin spells until the number of spells on your list matches the number in the Paladin Features table. The chosen spells must be of a level for which you have spell slots. For example, if you’re a level 5 Paladin, your list of prepared spells can include six Paladin spells of level 1 or 2 in any combination.
	
	If another Paladin feature gives you spells that you always have prepared, those spells don’t count against the number of spells you can prepare with this feature, but those spells otherwise count as Paladin spells for you.
	
	Changing Your Prepared Spells. Whenever you finish a Long Rest, you can replace one spell on your list with another Paladin spell for which you have spell slots.
	
	Spellcasting Ability. Charisma is your spellcasting ability for your Paladin spells.
	
	Spellcasting Focus. You can use a Holy Symbol as a Spellcasting Focus for your Paladin spells.
\end{DndComment}

\begin{DndComment}{Weapon Mastery}
	Your training with weapons allows you to use the mastery properties of two kinds of weapons of your choice with which you have proficiency.
	
	Whenever you finish a Long Rest, you can change the kinds of weapons you chose.
	
	\textbf{Quarterstaff (Topple).} Your training with weapons allows you to use the mastery property of Quarterstaffs:
	
	\textbf{Topple:} If you hit a creature with a Quarterstaff, you can force the creature to make a Constitution saving throw (DC 8 plus the ability modifier used to make the attack roll and your Proficiency Bonus). On a failed save, the creature has the Prone condition.
	
	\textbf{Longsword (Sap)} Your training with weapons allows you to use the mastery property of Longswords:
	
	\textbf{Sap:} If you hit a creature with a Longsword, that creature has Disadvantage on its next attack roll before the start of your next turn.
	
	\textbf{Topple (Quarterstaff):} 1 Action
	
	\textbf{Sap (Longsword):} 1 Action

\end{DndComment}


\begin{DndComment}{Channel Divinity}
	You can channel divine energy directly from the Outer Planes, using it to fuel magical effects. You start with one such effect: Divine Sense, which is described below. Other Paladin features give additional Channel Divinity effect options. Each time you use this class’s Channel Divinity, you choose which effect from this class to create.
	
	You can use this class’s Channel Divinity twice. You regain one of its expended uses when you finish a Short Rest, and you regain all expended uses when you finish a Long Rest. You gain an additional use when you reach Paladin level 11.
	
	If a Channel Divinity effect requires a saving throw, the DC equals the spell save DC from this class’s Spellcasting feature.
	
	\textbf{Divine Sense:} As a Bonus Action, you can open your awareness to detect Celestials, Fiends, and Undead. For the next 10 minutes or until you have the Incapacitated condition, you know the location of any creature of those types within 60 feet of yourself, and you know its creature type. Within the same radius, you also detect the presence of any place or object that has been consecrated or desecrated, as with the Hallow spell.
	
	
\end{DndComment}



\endgroup

